% ----------------------------------------------------------
% Capítulo de Resultados
% ----------------------------------------------------------
\chapter{Resultados}
\label{cap:resultados}
% ---
Este capítulo descreve os resultados finais do desenvolvimento da ferramenta proposta, abordando as funcionalidades implementadas, as limitações identificadas, bem como possíveis trabalhos futuros e expansões do projeto.

Foi desenvolvido um construtor gráfico de \ac{APN} baseado em grafos manipuláveis pelo usuário, além de um painel que ilustra de forma dinâmica a execução de uma determinada palavra e os múltiplos caminhos possíveis gerados pelo não-determinismo. Esse painel também destaca visualmente os caminhos que resultaram em aceitação. Ao submeter uma palavra à análise do autômato, o usuário pode selecionar entre três critérios de aceitação: pilha vazia (S), estado final (F) ou ambos (FS). O design da interface foi projetado priorizando a usabilidade e a agilidade nas interações em detrimento de configurações excessivamente detalhadas, focando na redução do número de cliques necessários para executar cada ação.

Os APNs modelados na ferramenta podem ser exportados em formato JSON e importados localmente para dar continuidade a estudos anteriores. Como o sistema foi concebido para fins essencialmente didáticos e de apoio gráfico, o foco recai sobre autômatos de pilha de pequeno porte, podendo apresentar degradação de desempenho conforme o grau de não-determinismo aumenta ou dependendo do comprimento da palavra analisada.

Dessa forma, a aplicação não foi projetada para testes massivos de modelos robustos, embora sua eficiência computacional possa ser refinada em etapas posteriores para cenários mais exigentes. O sistema também não mitiga de forma nativa a ocorrência de \textit{loops} infinitos no autômato; por essa razão, foi integrado um parâmetro configurável de limite máximo de iterações. Essa decisão justifica-se pelo fato de que a detecção e prevenção absoluta de ciclos infinitos em APNs constitui um problema complexo na literatura da área. Ademais, como os navegadores web modernos operam sua interface (\textit{frontend}) prioritariamente em uma única \textit{thread} e o projeto prescinde de um \textit{backend} — decisão estratégica para viabilizar a hospedagem gratuita e estável do sistema —, todo o processamento dos caminhos não-determinísticos é executado de forma síncrona no cliente.

O desenvolvimento do \textit{software} para abranger outras máquinas de estados chegou a ser iniciado, contudo, em razão do cronograma reduzido, optou-se pela delimitação do escopo aos APNs. A integração dessas demais máquinas permanece como proposta para trabalhos futuros. Por fim, ressalta-se que não foi conduzido um estudo empírico com estudantes para avaliar a real eficácia pedagógica da ferramenta no ensino da disciplina, consolidando a validação com usuários como uma das principais vertentes para a continuidade desta pesquisa.

